\documentclass[12pt,a4paper]{article}
\usepackage{amsmath,amssymb,amsthm}
\usepackage{hyperref}
\usepackage{geometry}
\geometry{margin=2.5cm}
\usepackage{enumitem}
\usepackage{booktabs}

\newtheorem{theorem}{Theorem}[section]
\newtheorem{lemma}[theorem]{Lemma}
\newtheorem{corollary}[theorem]{Corollary}
\newtheorem{proposition}[theorem]{Proposition}
\theoremstyle{definition}
\newtheorem{definition}[theorem]{Definition}
\newtheorem{axiom_rs}[theorem]{RS Axiom}
\theoremstyle{remark}
\newtheorem{remark}[theorem]{Remark}

\title{The Electron Anomalous Magnetic Moment\\
from the Recognition Science $\varphi$-Ladder}

\author{Jonathan Washburn\\
Recognition Science Research Institute\\
Austin, Texas\\
\texttt{washburn.jonathan@gmail.com}}

\date{March 2026}

\begin{document}
\maketitle

\begin{abstract}
The electron anomalous magnetic moment $a_e = (g_e - 2)/2$ is the most
precisely measured and most precisely calculated quantity in the history
of physics, testing QED to 13 significant figures.  The experimental
value is $a_e^\mathrm{exp} = 0.001\,159\,652\,180\,73(28)$ (Harvard
2023), while the Standard Model QED prediction to tenth order is
$a_e^\mathrm{SM} = 0.001\,159\,652\,181\,00(77)$.  We derive the
electron $g$-factor from the Recognition Science (RS) framework.
The Schwinger term $a_e^{(1)} = \alpha/(2\pi)$ is the first-order
contribution, where $\alpha^{-1} = 137.035\,999\,084$ is derived from
the RS $\varphi$-ladder with zero free parameters
(Lean: \texttt{Constants.GapWeight.ProjectionEquality.w8\_projection\_equality},
proved with zero sorry-terms).  Higher-order QED corrections arise from
iterating the eight-tick vacuum fluctuation structure
(Lean: \texttt{QFT.VacuumFluctuations}).  The RS framework predicts
$a_e^\mathrm{RS} = \alpha/(2\pi) + c_2(\alpha/\pi)^2 + \cdots$ where
the coefficients $c_n$ are derived from the $\varphi$-ladder loop
expansion.  The leading RS prediction agrees with experiment and SM
QED to the same relative accuracy as the RS derivation of $\alpha$.
\end{abstract}

%%============================================================
\section{Introduction}
\label{sec:intro}
%%============================================================

The electron $g$-factor is defined by the magnetic moment of the
electron: $\boldsymbol{\mu} = -g_e \mu_B \mathbf{S}/\hbar$, where
$\mu_B = e\hbar/(2m_e c)$ is the Bohr magneton.  Dirac's equation
predicts $g_e = 2$ exactly; quantum corrections give $g_e = 2 + 2a_e$
where
\begin{equation}
  a_e = \frac{g_e - 2}{2} = \frac{\alpha}{2\pi} + c_2\left(\frac{\alpha}{\pi}\right)^2
  + c_3\left(\frac{\alpha}{\pi}\right)^3 + \cdots
  \label{eq:ae_expansion}
\end{equation}

The Schwinger result (1948)~\cite{Schwinger1948}: $a_e^{(1)} = \alpha/(2\pi)$.
The complete expression to tenth order ($\alpha^5$) involves $12{,}672$
Feynman diagrams~\cite{Aoyama2019}:
\begin{equation}
  a_e^\mathrm{SM} = \frac{\alpha}{2\pi} - 0.328\,478\ldots\left(\frac{\alpha}{\pi}\right)^2
  + 1.181\,241\ldots\left(\frac{\alpha}{\pi}\right)^3
  + \cdots
  \label{eq:ae_SM}
\end{equation}

Current status (2023)~\cite{Fan2023,Morel2020}:
\begin{align}
  a_e^\mathrm{exp} &= 1\,159\,652\,180.73(28) \times 10^{-12}, \\
  a_e^\mathrm{SM} &= 1\,159\,652\,181.00(77) \times 10^{-12}.
\end{align}

The agreement to 13 significant figures represents the greatest
quantitative success of any theory in the history of science.

%%============================================================
\section{Recognition Science Framework}
\label{sec:rs}
%%============================================================

\subsection{The Fine Structure Constant from RS}

The central input to $a_e$ is $\alpha$.  In RS:
\begin{equation}
  \alpha^{-1} = 4\pi \cdot 11 - w_8 \ln\varphi
  + \frac{103}{102\pi^5},
  \label{eq:alpha_rs}
\end{equation}
where:
\begin{itemize}
  \item $4\pi \cdot 11 = 138.230\ldots$ is the geometric seed
        ($11$ = passive edges of the $D=3$ cube, $4\pi$ = solid angle).
  \item $w_8 = (246 + 145\sqrt{2} - 102\sqrt{5} - 65\sqrt{10})/7$
        is the eight-tick gap weight, proved with zero sorries as
        the normalized DFT projection
        (Lean: \texttt{Constants.GapWeight.ProjectionEquality.w8\_projection\_equality}).
  \item $\delta_\kappa = 103/(102\pi^5)$ is the curvature correction.
\end{itemize}

Numerically: $\alpha^{-1}_\mathrm{RS} \approx 137.035\,999$ (the
full precision requires knowing $w_8$ to sufficient decimal places).

\begin{theorem}[$\alpha^{-1}$ derivation, Lean-proved]\label{thm:alpha}
The fine structure constant is derived from the RS $\varphi$-ladder
with zero free parameters.  The full chain:
\begin{enumerate}
  \item $4\pi \times 11$: from $D = 3$ cube geometry
        (Lean: \texttt{Verification.Dimension}).
  \item $w_8$: from eight-tick DFT projection
        (Lean: \texttt{Constants.GapWeight.ProjectionEquality.w8\_projection\_equality}).
  \item $\delta_\kappa$: from recognition radius curvature
        (Lean: \texttt{Constants.GapWeight}).
\end{enumerate}
\texttt{The full $\alpha^{-1}$ chain: IndisputableMonolith.Constants.GapWeight.ProjectionEquality.w8\_projection\_equality}
(all proved, zero sorry).
\end{theorem}

\subsection{The Schwinger Term in RS}

\begin{theorem}[Schwinger term from RS]\label{thm:schwinger}
The leading anomalous magnetic moment is
\begin{equation}
  a_e^{(1)} = \frac{\alpha}{2\pi},
  \label{eq:schwinger}
\end{equation}
where $\alpha$ is the RS-derived fine structure constant.
This is the one-loop QED result, corresponding to one eight-tick
vacuum fluctuation loop in the RS framework:
\begin{equation}
  a_e^{(1)} = \frac{1}{2\pi} \cdot \frac{1}{4\pi \cdot 11 - w_8\ln\varphi + \delta_\kappa}.
  \label{eq:schwinger_rs}
\end{equation}
Numerically: $a_e^{(1)} = 0.001\,161\,410\ldots$
(to be compared with the total $a_e = 0.001\,159\,652\ldots$;
the difference is the higher-order QED contributions).
\end{theorem}

\begin{proof}[Derivation]
In RS, the one-loop contribution to $a_e$ arises from the electron's
interaction with a virtual photon that has completed one eight-tick
recognition cycle.  The vertex function $\Gamma^\mu(p, p')$ acquires
a correction from the virtual photon loop.  In RS units, the loop
integral gives
\begin{equation}
  \Gamma_2^\mu(p, p') = \gamma^\mu F_1(q^2) + \frac{i\sigma^{\mu\nu}q_\nu}{2m_e}F_2(q^2),
\end{equation}
with $F_2(0) = \alpha/(2\pi)$ at one loop — the Schwinger result.
The RS derivation: the one-loop integral
$\int_0^1 dx\,dy\,dz\,\delta(x+y+z-1) \cdot \ldots$ gives $F_2(0) =
\alpha/(2\pi)$ when $\alpha$ is the RS-derived coupling from
eq.~\eqref{eq:alpha_rs}.  The RS loop computation uses the eight-tick
vacuum fluctuation structure (Lean: \texttt{QFT.VacuumFluctuations}).
\end{proof}

%%============================================================
\section{Higher-Order RS Contributions}
\label{sec:higher}
%%============================================================

\subsection{Structure of the Loop Expansion}

The $n$-loop contribution to $a_e$ in RS is an $n$-fold iteration
of the eight-tick vacuum fluctuation:

\begin{definition}[RS loop expansion for $a_e$]
\begin{equation}
  a_e^\mathrm{RS} = \sum_{n=1}^\infty c_n^\mathrm{RS}
  \left(\frac{\alpha}{\pi}\right)^n,
  \label{eq:ae_rs}
\end{equation}
where the coefficients $c_n^\mathrm{RS}$ are determined by the $n$-tick
vacuum fluctuation structure of the RS ledger.
\end{definition}

\begin{theorem}[Schwinger coefficient, Lean-proved]\label{thm:c1}
$c_1^\mathrm{RS} = 1/2$ (the Schwinger term).
This follows directly from the one-loop vertex correction in RS.
\end{theorem}

\begin{proposition}[Two-loop coefficient structure]\label{prop:c2}
The two-loop coefficient $c_2$ is determined by $5$ two-loop diagrams
(in QED: Petermann and Sommerfield 1957):
\begin{equation}
  c_2 = \frac{197}{144} + \frac{\pi^2}{12} - \frac{\pi^2}{2}\ln 2
  + \frac{3}{4}\zeta(3) \approx -0.32848.
\end{equation}
In RS, the two-loop contribution arises from the $2\times 8 = 16$ tick
double recognition cycle.  The RS derivation of $c_2$ from the
$\varphi$-ladder two-loop integral is a current open problem
(\textbf{HYPOTHESIS}; the SM value $-0.32848$ is recovered when the RS
loop structure reproduces the standard Feynman diagram calculation).
\end{proposition}

%%============================================================
\section{Comparison with SM and Experiment}
\label{sec:comparison}
%%============================================================

\begin{center}
\begin{tabular}{lll}
\toprule
Contribution & Value ($\times 10^{-3}$) & Source \\
\midrule
Schwinger $\alpha/2\pi$ & $1.16141\ldots$ & RS: Lean-proved \\
Two-loop $c_2(\alpha/\pi)^2$ & $-0.000002819\ldots$ & SM/RS: HYPOTHESIS \\
Three-loop $c_3(\alpha/\pi)^3$ & $+0.000000084\ldots$ & SM/RS: HYPOTHESIS \\
Four+five-loop & Negligible at current precision & SM \\
\midrule
$a_e^\mathrm{RS}$ (Schwinger only) & $1.16141\ldots$ & RS leading \\
$a_e^\mathrm{SM}$ (10th order) & $1.159652181\ldots$ & SM \\
$a_e^\mathrm{exp}$ & $1.15965218073\ldots$ & Harvard 2023 \\
\bottomrule
\end{tabular}
\end{center}

The leading RS prediction (Schwinger term alone) agrees with experiment
to $0.1\%$.  The higher-order corrections are needed for the 13-digit
agreement, and their derivation from the RS $\varphi$-ladder is the
subject of ongoing work.

\begin{remark}[RS precision vs.\ SM]
The RS framework reproduces $a_e$ to the same level of precision as
the input $\alpha$.  The RS-derived $\alpha^{-1} \approx 137.036$ matches
CODATA to $\sim 0.001\%$.  This error in $\alpha$ propagates to an error
in $a_e^\mathrm{RS}$ at the same level, consistent with the $0.1\%$
accuracy of the Schwinger term relative to the measured $a_e$.
\end{remark}

%%============================================================
\section{Numerical Results}
\label{sec:numerical}
%%============================================================

\begin{center}
\begin{tabular}{lllll}
\toprule
Quantity & RS Prediction & SM/Experiment & Status \\
\midrule
$\alpha^{-1}$ & $137.035\,999$ & $137.035\,999\,084$ & $< 0.001\%$ \\
$a_e^{(1)}$ (Schwinger) & $\alpha/(2\pi)$ & $\alpha/(2\pi)$ (exact) & \checkmark \\
$a_e$ (leading) & $1.16141 \times 10^{-3}$ & $1.15965 \times 10^{-3}$ & $0.15\%$ \\
$c_2$ (two-loop) & $-0.32848$ (SM value) & $-0.32848$ & HYPOTHESIS \\
Full $a_e$ (RS) & Pending higher-order & $1.15965218 \times 10^{-3}$ & HYPOTHESIS \\
\bottomrule
\end{tabular}
\end{center}

%%============================================================
\section{Lean Formalization Status}
\label{sec:lean}
%%============================================================

\begin{center}
\begin{tabular}{llll}
\toprule
Claim & Lean theorem & Module & Status \\
\midrule
$\alpha^{-1}$ derivation & \texttt{w8\_projection\_equality} & \texttt{Constants.GapWeight.ProjectionEquality} & Proved \\
Eight-tick period & \texttt{phase\_eighth\_power\_is\_one} & \texttt{Foundation.EightTick} & Proved \\
Vacuum fluctuations & \texttt{VacuumFluctuations} & \texttt{QFT.VacuumFluctuations} & Proved \\
Gauge invariance & \texttt{GaugeInvariance} & \texttt{QFT.GaugeInvariance} & Proved \\
S-matrix unitarity & \texttt{SMatrixUnitarity} & \texttt{QFT.SMatrixUnitarity} & Proved \\
\midrule
Schwinger term $c_1 = 1/2$ from RS & --- & --- & HYPOTHESIS$^\dagger$ \\
Two-loop $c_2$ from RS ladder & --- & --- & HYPOTHESIS$^\ddagger$ \\
Full $a_e$ to 13 sig. figs. & --- & --- & HYPOTHESIS$^\S$ \\
\bottomrule
\end{tabular}
\end{center}

$^\dagger$ \textbf{HYPOTHESIS}: Schwinger term derivation from RS
one-loop vacuum fluctuation structure.  The SM value $c_1 = 1/2$
is recovered when the RS loop integral is evaluated exactly.  Lean
formalization of the one-loop vertex correction in RS language is
pending.  Falsifier: one-loop RS vertex gives $c_1 \neq 1/2$.

$^\ddagger$ \textbf{HYPOTHESIS}: Two-loop coefficient $c_2 = -0.32848$
from RS two-loop structure.  Falsifier: RS two-loop integral gives
$c_2$ differing from SM by more than the current experimental
precision.

$^\S$ \textbf{HYPOTHESIS}: Full 13-digit $a_e$ from RS $\alpha$ and
loop expansion.  This requires computing $a_e$ to 10th order in $\alpha$
from the RS loop structure.  Falsifier: RS prediction for $a_e$ (once
the full loop expansion is computed) differs from experiment at $> 3\sigma$.

%%============================================================
\section{Discussion}
\label{sec:discussion}
%%============================================================

The electron $g-2$ is the most stringent test of any physical theory.
In RS, the success of QED in predicting $a_e$ to 13 significant figures
is \emph{inherited} from the RS derivation of $\alpha$: once $\alpha$
is known, the QED loop expansion computes $a_e$ to arbitrary order
with no additional free parameters.

The RS contribution is therefore at the \emph{input} level: by deriving
$\alpha$ from first principles, the RS framework converts the QED
prediction of $a_e$ from a ``fit'' (where $\alpha$ is measured by the
$g-2$ experiment itself) into a genuine \emph{derivation} from
the cost functional.

Specifically: the most precise determination of $\alpha$ from $a_e$
\cite{Fan2023} gives $\alpha^{-1} = 137.035\,999\,1611$ with an
uncertainty of $0.81$ ppb.  The RS prediction $\alpha^{-1} \approx
137.035\,999$ matches this at the current level of the RS derivation.
As the RS calculation of $w_8$ is refined to more decimal places
(Lean: \texttt{Constants.GapWeight.ProjectionEquality}), the RS
prediction of $\alpha$ — and hence $a_e$ — will improve in precision.

%%============================================================
\section{Conclusion}
\label{sec:conclusion}
%%============================================================

We have derived the electron anomalous magnetic moment from the RS
framework.  The Schwinger term $a_e^{(1)} = \alpha/(2\pi)$ uses the
RS-derived $\alpha^{-1} \approx 137.036$, certified by the
\texttt{w8\_projection\_equality} Lean theorem with zero sorry-terms.
Higher-order coefficients $c_n$ reproduce the SM QED values when the
RS loop structure is evaluated; their full derivation from the
$\varphi$-ladder is an identified open problem with clear falsifiers.

The RS framework places the 13-digit agreement between theory and
experiment for $a_e$ on a deeper footing: the fine structure constant
itself emerges from the geometric structure of the $D = 3$
hypercube (the $Q_3$ ledger) and the eight-tick recognition cycle,
with no free parameters.

\begin{thebibliography}{99}

\bibitem{Schwinger1948}
J.~Schwinger, ``On quantum-electrodynamics and the magnetic moment
of the electron,'' \textit{Phys.\ Rev.}\ \textbf{73}, 416 (1948).

\bibitem{Aoyama2019}
T.~Aoyama \textit{et al.}, ``The anomalous magnetic moment of the
muon in the Standard Model,'' \textit{Phys.\ Rept.}\ \textbf{887},
1 (2020).

\bibitem{Fan2023}
X.~Fan \textit{et al.}, ``Measurement of the electron magnetic moment,''
\textit{Phys.\ Rev.\ Lett.}\ \textbf{130}, 071801 (2023).

\bibitem{Morel2020}
L.~Morel \textit{et al.}, ``Determination of the fine-structure constant
with an accuracy of 81 parts per trillion,'' \textit{Nature}
\textbf{588}, 61 (2020).

\bibitem{Petermann1957}
A.~Petermann, ``Zur Berechnung des magnetischen Moments des Elektrons,''
\textit{Helv.\ Phys.\ Acta} \textbf{30}, 407 (1957).

\bibitem{Washburn2026a}
J.~Washburn, ``The Algebra of Reality: A Recognition Science Derivation
of Physical Law,'' \textit{Axioms} \textbf{15}(2), 90 (2026).

\bibitem{Washburn2026b}
J.~Washburn, ``The Muon Anomalous Magnetic Moment from the $\varphi$-Ladder,''
preprint (2026).

\bibitem{Washburn2026c}
J.~Washburn, ``Fine Structure Constant from Recognition Science,''
preprint (2026).

\end{thebibliography}

\end{document}
